\scanpage{159}
% L1–3
trouve une structure d'extension
\diagram{\[
\begin{tikzcd}[ampersand replacement=\&]
1 \arrow[r,dash] \& G_0 \arrow[r,dash] \& G \arrow[r,"\chi'"] \& Z \arrow[r,dash] \& 1\ . \\
 \& \& Z \arrow[u,hook] \arrow[ur,"{\lambda\mapsto\lambda^2}"'] \& \&
\end{tikzcd}\tag*{(81)}
\]}
% slanted
\margin{On va supposer par la suite que l'on n'est pas dans le cas octaédral, avec $Z=\{1,\unsure{=|-|\varepsilon}\}$.}

% L4–9
On peut introduire une notion d'\uline{équivalence} de deux structures de 2-polyèdres \deleted{réguliers} triangulés \inserted{sur le $=$ ens. de sommets,} munis d'un groupe $Z$ \deleted{\unsure{opérant}} de permutations de $S$ (jouant le rôle d'``\uline{homothéties}'' -- on parlera aussi de \unsure{g|2} polyèdres à opérateurs.
% archive sheet number: 158
